\section{Introduction}
\label{sec:intro}

Vehicles in a V2X network broadcast their position, speed, and acceleration in periodic safety messages, and
neighboring vehicles act on that content. Cryptographic credentials show that a sender is enrolled; they do not show
that the content is true. Misbehavior detection addresses this second problem by checking the plausibility of what
a vehicle receives~\cite{heijden2018veremi,kamel2020veremi}. A Sybil attacker is a hard case: one vehicle sends
under several pseudonyms, so that a receiver sees several apparent vehicles, for example a traffic jam that does
not exist~\cite{kamel2019misbehavior}.

Most learning-based detectors for this problem are developed on VeReMi~\cite{heijden2018veremi} and
VeReMi-Extension~\cite{kamel2020veremi}, and most are supervised~\cite{boualouache2023survey}. The reported
numbers are high. Supervised sequence classifiers report accuracy or F1 between 0.95 and
0.998~\cite{alladi2023deep,hamhoum2024mistralbsm,youness2025vemisnet}. Reconstruction-based detectors trained on
benign data only, which need no attack labels and are therefore attractive for attacks not seen in training, report
0.93 to 0.99~\cite{alladi2021deepadv,alladi2021deep,kalogiannis2026pampos,zhong2023sybil}, with the four Sybil
attacks of VeReMi-Extension among the best-detected
classes~\cite{alladi2021deep,kalogiannis2026pampos}.

The protocols behind these numbers vary, and several of their elements are often not stated. Neither dataset paper
defines a train/test split~\cite{heijden2018veremi,kamel2020veremi}. Where a split is described, it is usually
drawn over windows or sequences~\cite{franceschini2025line,youness2025vemisnet}; among the papers we read, only one
states a split by sender~\cite{kalogiannis2026pampos}. Sequences are formed per sender identifier in some
studies~\cite{hamhoum2024mistralbsm} and require a minimum number of messages in
others~\cite{kalogiannis2026pampos}. Decision thresholds are taken from benign validation data in some
studies~\cite{kalogiannis2026pampos} and selected on labeled test data in others~\cite{alladi2021deep}. The term
``unsupervised'' almost always means training on benign data selected with the ground-truth labels. To our
knowledge, no study compares several detectors and several training-data assumptions on the VeReMi-Extension Sybil
attacks under one protocol, and none reports how much of a detector's score a trivial property of the input, such
as the number of messages in a window, already explains. The machine-learning and security literature gives reasons
to ask both questions~\cite{arp2022dos,geirhos2020shortcut,kapoor2023leakage,wu2021current}.

The second question matters for Sybil attacks in particular. In three of the four attacks (DataReplaySybil,
DoSRandomSybil, DoSDisruptiveSybil) the attacker changes pseudonym every one or two messages. A receiver that
groups messages by pseudonym, the only identifier it can observe, therefore sees these attacks as very short
sequences: in the data used here, 95\% to 99.7\% of their windows contain one to three messages, against about
14\% of benign windows. Sequence length alone then separates the classes (AUROC \prelim{0.95--0.97} in our earlier
studies), whatever a model learns about motion.

This paper asks: \emph{with vehicles held out, how well can each Sybil attack of VeReMi-Extension be separated from
benign traffic, as a function of what the training data is assumed to be?} We compare three assumptions (benign
traffic only, an unlabeled mix of benign and attack traffic, and labeled data) across six unsupervised detector
families and two supervised models, on one input representation, one split by physical vehicle, and one test set.
The protocol, including the rules for reading the results, was fixed before the runs (Section~\ref{sec:protocol}).
The claim under test is that the training-data assumption changes detection more than the choice of detector does,
and that for the three pseudonym-rotating attacks window length explains most of any detector's score.
\pending{Whether the results support this claim is reported in Section~\ref{sec:results}; this paragraph is
completed when the runs are finished.}

The contributions are as follows.
\begin{itemize}
  \item A vehicle-disjoint evaluation protocol for the four Sybil attacks of VeReMi-Extension, with receiver-side
  pseudonym windows, thresholds from validation data only, five seeds, and confidence intervals from a bootstrap
  over vehicles (Section~\ref{sec:protocol}).
  \item A comparison, under this protocol, of convolutional and LSTM autoencoders, PCA reconstruction, Isolation
  Forest, one-class SVM, and $k$-nearest-neighbor distance under benign-only and unlabeled-mix training, and of two
  supervised models (Sections~\ref{sec:detectors} and~\ref{sec:results}). To our knowledge, no such comparison has
  been reported for these attacks.
  \item A length-only reference score and a length-matched AUROC for every detector, which separate what a detector
  adds from what the number of messages in a window already gives for the pseudonym-rotating attacks. In our
  earlier studies, matching on length turned an apparent DataReplay AUROC of \prelim{0.69} into chance
  (\prelim{0.49}) for a masked autoencoder (Section~\ref{sec:disc-length}).
  \item A measurement, on Sybil attacks, of failure modes that are known in general anomaly detection: training on
  data in which the attack is frequent, and reconstruction of replayed genuine trajectories
  (Section~\ref{sec:discussion}). \pending{Wording to be adjusted to the outcome.}
\end{itemize}

We do not attribute the difference between our numbers and published ones to a single cause: the published values
are accuracy, F1, or recall at a chosen threshold on other sequence units, and the protocol here varies neither the
split nor the sequence unit.

Section~\ref{sec:background} describes the dataset, the attacks, and related work. Section~\ref{sec:protocol}
gives the protocol and Section~\ref{sec:detectors} the detectors. Section~\ref{sec:results} reports the results,
and Section~\ref{sec:discussion} discusses them and their limitations.
